\documentclass[11pt]{article}
\usepackage{amsmath,amssymb}
\usepackage{tikz}
\begin{document}
\section{Introduction}\label{sec:intro}
Every writer knows the rhythm.

\begin{figure}[h!]
\centering
\newcommand\petals{5}
\begin{tikzpicture}[scale=.6]
  \foreach \i in {1,...,\petals}
    \draw[fill=red!15, rotate=\i*360/\petals] (0,0) .. controls (.7,.45) and (1.7,.35) .. (2,0) .. controls (1.7,-.35) and (.7,-.45) .. (0,0);
\end{tikzpicture}
\caption{A figure.}\label{fig:live}
\end{figure}

Section~\ref{sec:method} describes the method and Section~\ref{sec:results} the results~\cite{knuth84}.

\newpage
\section{Method}\label{sec:method}
A document is a program.
\[ x = 1 \]

\newpage
\section{Results}\label{sec:results}
It works.

\begin{thebibliography}{9}
\bibitem{knuth84} D.~E. Knuth. \emph{The \TeX book}. 1984.
\end{thebibliography}
\end{document}
